\section[Global VWF compression]{Global signed VWF compression with an absolute error budget}
\label{sec:op3-global-vwf-compression}

\paragraph{Status and input contract.}
These are \statusproved{} drafts for a supplied scalar function, not a
recursive diffusion solver. The input is a vertex weighting function on
$[L,\infty)$, $L\leq0$, given by its sorted curvature-drop events:
\begin{equation}
 F(x)=c+ax+\sum_{j=1}^k q_j B_{s_j}(x),\qquad
 c\leq0,\quad q_j>0,\quad L<s_1<\cdots<s_k,
 \label{eq:op3-vwf-atoms}
\end{equation}
where
\[
 B_s(x)=
 \begin{cases}
  (s-x)_+^2/2,&s\leq0,\\
  x^2/2-(x-s)_+^2/2,&s\geq0.
 \end{cases}
\]
The formulas agree at $s=0$. Each atom is convex in $x$, has value and
derivative zero at zero, and has second derivative $\one_{x<s}$ off its
breakpoint. Thus $c=F(0)$, $a=F'(0)$, and $C=\sum_jq_j$ bounds the initial
curvature. A supplied piecewise quadratic VWF has this representation by
listing its nonzero curvature drops; the constant final derivative makes
the remaining quadratic coefficient zero. Conversion costs $O(1+k)$.

Chen--Peng--Wang's geometric compression uses the same signed atoms and
rounds breakpoints away from zero; see Lemma 7.7 and the proof of Lemma 6.4,
PDF pp.\ 37--39 \citep{chen2021diffusion}. Their size bound uses the
encountered-number assumption through Fact 7.2. The following regularization
removes dependence on the smallest nonzero breakpoint for this operation.
It does not bound every number encountered by the source algorithm.

\begin{lemma}[Tiny-breakpoint moment regularization; \statusproved]
\label{lem:op3-vwf-moment-regularization}
Fix $\tau>0$. For every event with $r=|s|<\tau$, replace $qB_s$ by
\[
 q(1-r/\tau)B_0+q(r/\tau)B_{\operatorname{sgn}(s)\tau}.
\]
For $s=0$ the second term is omitted. For each negative such event also
subtract the constant $A_s=qr(\tau-r)/2$. Leave all other events unchanged.
Call the result $F_{\rm reg}$ and set
\begin{equation}
 \xi=\sum_{|s_j|<\tau}\frac{q_j|s_j|(\tau-|s_j|)}2
 \leq \frac{C\tau^2}{8}.
 \label{eq:op3-global-regularization-error}
\end{equation}
Then, on the entire original domain,
\[
 0\leq F(x)-F_{\rm reg}(x)\leq\xi.
\]
The result is a VWF, has the same derivative at zero and the same final
derivative as $F$, and has no greater initial curvature.
\end{lemma}
\begin{proof}
For fixed $x$, the function $s\mapsto B_s(x)$ on $[0,\tau]$ is concave,
with second derivative between $-1$ and zero wherever it exists. Its chord
lies below it by at most $r(\tau-r)/2$. This last bound follows by
subtracting the chord and comparing with $s(\tau-s)/2$: their difference
is convex and zero at both endpoints. On $[-\tau,0]$ the same function
is convex with second derivative between zero and one. Its chord lies
above it by at most $r(\tau-r)/2$. Subtracting this upper bound therefore
makes the negative-event replacement a lower approximation as well.
These statements also hold at the piecewise quadratic joins by continuity.
Sum the individual nonnegative errors and use $r(\tau-r)\leq\tau^2/4$.

All replacement curvature weights are nonnegative and have the same total
as their input event. The new value at zero is
$c-\sum_{-\tau<s_j<0}A_{s_j}\leq0$, and the derivative there remains $a$.
The final derivative is $a+\sum_{s_j>0}q_js_j$. Every positive replacement
preserves its contribution because $(qr/\tau)\tau=qr$; negative and zero
atoms contribute zero. An event at or below $L$ is identically zero on
the domain and may be omitted. Such omissions can only decrease the
listed initial curvature.
\end{proof}

\begin{theorem}[Global dyadic compression; \statusproved]
\label{thm:op3-global-vwf-compression}
Let $R=\max\{\tau,\max_j|s_j|\}$. After the regularization above, round
every nonzero event $s$ away from zero to the smallest signed dyadic grid
point $t$ of magnitude at least $|s|$, on the grid
$\tau,2\tau,4\tau,\ldots$. Replace its weight $q$ by $q|s|/|t|$.
Keep zero events explicitly without taking a ratio. The resulting VWF
$\widehat F$ satisfies, for every $x\geq L$,
\begin{equation}
 F(x)\geq\widehat F(x)\geq2F(x/2)-2\xi.
 \label{eq:op3-global-vwf-compression}
\end{equation}
It preserves the final derivative, has initial curvature at most $C$, and
has at most $2(1+\lceil\log_2(R/\tau)\rceil)+1$ events. Every nonzero
output breakpoint has magnitude between $\tau$ and $2R$.
Construction costs $O(1+k+\log(R/\tau))$ exact-real word operations and
$O(1+\log(R/\tau))$ additional allocated words, including the output.
The input event storage is supplied and is not counted as new allocation.
\end{theorem}
\begin{proof}
For any signed $s$ and any $b>0$, direct substitution gives
$B_{bs}(x)=b^2B_s(x/b)$. Since $B_s$ is convex and $B_s(0)=0$,
$bB_s(x/b)$ is nonincreasing in $b>0$: for $0<b_1\leq b_2$, apply
convexity to $x/b_2=(b_1/b_2)(x/b_1)+(1-b_1/b_2)0$.
For a rounded event, $t=bs$ with $1\leq b\leq2$, hence
\[
 2qB_s(x/2)\leq (q/b)B_{bs}(x)\leq qB_s(x).
\]
This proof applies to positive and negative $x,s$. A zero event is kept
exactly and obeys the same inequalities by convexity. Summing and using
$F_{\rm reg}(0)\leq0$ gives
$2F_{\rm reg}(x/2)\leq\widehat F(x)\leq F_{\rm reg}(x)$.
Here $x/2\geq L$, since $L\leq0$. The preceding lemma now gives
\eqref{eq:op3-global-vwf-compression}. Rounding preserves $qs$ for positive
events and never increases a curvature weight. The constant and linear
terms are retained, proving the derivative claims.

One scan finds $R$. Repeated doubling builds the finite grid; no constant-time
real logarithm oracle is assumed. A second scan of the sorted events fills
positive, negative and zero bins. A negative-side pointer moves monotonically
toward zero, and a positive-side pointer moves monotonically away from zero.
Tiny events enter just the three bins $-\tau,0,\tau$, including the constant
shift. Each pointer visits the grid at most once. Finally stream the bins
in sorted order, omitting zero weights and events at or below $L$.
All scans, bins, output records and allocations fit the stated bounds;
neither comparison sorting nor an uncharged hash table is needed.
\end{proof}

\paragraph{Choice of the absolute budget.}
For a requested $\xi_0>0$ and $C>0$, the conservative word-computable choice
$\tau=\min\{1,\xi_0/C\}$ ensures $C\tau^2/8\leq\xi_0$.
If $C=0$, return the affine input exactly. This does not require extracting
a square root. The size bound then depends logarithmically on the largest
split, curvature upper bound and requested error. Tiny curvature weights
can still remain: $q=1$, $s=2^{-2048}$, $\tau=2^{-10}$ creates a weight
$2^{-2038}$ at $\tau$. Thus the theorem is not a proof of the source's
Assumption 3.15, even though the tiny split itself has disappeared.

\begin{corollary}[One compressed graph-objective embedding; \statusproved]
\label{cor:op3-global-vwf-energy-embedding}
Let $E(x)=x^\top\mathcal Lx/2+\sum_iF_i(x_i)$ on
$x_i\geq L_i$, where $\mathcal L$ is a graph Laplacian and $L_i\leq0$.
Compress each vertex function as above and let $\Xi=\sum_i\xi_i$.
For the resulting objective $\widehat E$,
\[
 E(x)\geq\widehat E(x)\geq2E(x/2)-2\Xi.
\]
Assume the infima are finite, $E^*<0$, $\beta\geq1$, and a feasible
candidate obeys $\widehat E(w)\leq\widehat E^*/\beta$.
Then $E(w/2)\leq E^*/(2\beta)+\Xi$.
If additionally $\Xi\leq|E^*|/(4\beta)$, it is a $4\beta$-approximation
in the negative-energy convention: $E(w/2)\leq E^*/(4\beta)$.
\end{corollary}
\begin{proof}
The upper bound follows term by term. The lower bound has an additional
nonnegative quadratic $x^\top\mathcal Lx/4$, besides the scalar lower
bounds. Since $\widehat E^*\leq E^*$, substitution of the candidate gives
the claim. Feasibility is preserved by division by two.
\end{proof}

\paragraph{Weaker bounded-domain construction retained for comparison.}
On $[\max\{L,-R\},R]$, dropping events at or below $-R$ and clamping
events above $R$ to $R$ changes no atom value. Moving every $|s|<\tau$
to zero changes its atom by at most $R|s|$: this follows directly from
the two formulas in \eqref{eq:op3-vwf-atoms}. Put
$\xi_R=R\sum_{|s|<\tau}q|s|$, and call the snapped function $G$, retaining
the original affine terms. Thus $|G-F|\leq\xi_R$ on this interval.
Its dyadic compression $H$ obeys $G\geq H\geq2G(x/2)$ by the perspective
argument above. Return $\widehat F_R=H-\xi_R$. Then
$F\geq\widehat F_R\geq2F(x/2)-3\xi_R$ on the bounded domain.
We use the stronger global theorem for subsequent arguments.
Deleting a tiny negative event outright is invalid: at $x=-1$,
$B_{-\tau/2}(x)=(1-\tau/2)^2/2$, which does not tend to zero with $\tau$.

\paragraph{Audit and remaining obligations.}
\texttt{global\_vwf\_compression.py} audits 1,212 supplied functions in
exact rational arithmetic, including splits of both signs down to
$2^{-2048}$. Independent polynomial overlays certify 23,000 complete
intervals, 87,152 finite-interval inequalities and 4,848 infinite-ray
inequalities; all 1,212 final derivatives agree exactly. All grid scans,
pointer moves and output allocations are charged. The separate bounded
baseline audits 1,206 functions and 18,616 complete finite intervals.
These are implementations of scalar compression, not of the source solver.
Recursive upper-split and curvature bounds, additive budgets through all
refinement calls, and local graph discovery with cumulative paid work
remain \statusopen. General OP3 is unchanged.
